% !TeX root = ../thesis.tex

\chapter{制造、检验与压力试验}

\section{制造技术要求}

\subsection{钢材进场检验}

Q345R 钢板进场后逐张进行外观检查（不得有裂纹、气泡、结疤、折叠和夹杂等缺陷）和尺寸检查。钢板按批号核验质量证明书，保证化学成分和力学性能符合 GB/T 713.1—2023 的规定\cite{gbt713-1-2023}。按 NB/T 47013.3—2015 逐张进行超声检测，合格级别不低于 Ⅰ 级\cite{nbt47013-2015}。20 号钢无缝钢管按 GB/T 6479—2023\cite{gbt6479-2023}（或 GB/T 8163—2018\cite{gbt8163-2018}）验收。

\subsection{成形与组装}

筒体采用冷卷或温卷成形。卷制前对钢板边缘进行预弯或采用带压边圈的卷板机。成形后的筒节椭圆度（最大内径与最小内径之差）不大于名义内径的 $1\%$（$18$ mm），棱角度（E）不大于 $(\delta_n/10 + 2)$ mm 且 $\le 5$ mm\cite{hgt20584-2020}。封头采用整体冲压或旋压成形，成形后封头实测最小厚度不小于设计有效厚度 $10.7$ mm。

筒节纵缝、筒体与封头环缝均采用双面埋弧焊（SAW），全焊透。坡口形式为对称双 V 形或 U 形。焊接材料按等强度匹配原则选取：焊丝选用 H10Mn2，焊剂选用 SJ101（碱性烧结焊剂）。所有 A、B 类焊接接头为全焊透对接型式。

\subsection{焊后热处理}

根据 GB/T 150.4—2024 第 4.6.1 条\cite{gbt150-2024}，Q345R 制容器仅在厚度 $>38$ mm，或厚度 $>34$ mm 且设计温度低于 $0$ ℃，或有应力腐蚀要求时，方强制进行焊后热处理。本设备壁厚仅 $14$ mm，介质为无应力腐蚀的冷却水，不满足强制热处理条件，不进行整体焊后热处理。考虑冷卷成形残余应力，可仅对纵缝进行局部消除应力处理。

\section{无损检测}

按 GB/T 150.4—2024 和 NB/T 47013—2015\cite{gbt150-2024,nbt47013-2015} 的规定执行如下检测：\\[4pt]
\textbf{射线检测（RT）}。A、B 类对接焊缝（纵缝和环缝）进行局部射线检测，检测比例不少于各焊缝长度的 $20\%$ 且不小于 $250$ mm。透照技术等级不低于 AB 级，像质计灵敏度按 NB/T 47013.2—2015 的相应要求，合格级别不低于 Ⅲ 级。T 字接头交叉部位必须纳入检测范围。\\[4pt]
\textbf{超声检测（UT）}。对因结构限制无法实施射线检测的对接焊缝部位，可以超声检测替代，检测比例同射线检测。合格级别不低于 NB/T 47013.3—2015 的 Ⅰ 级。\\[4pt]
\textbf{渗透检测（PT）}。所有 C、D 类角焊缝（接管与壳体角焊缝、法兰与接管角焊缝、补强圈与壳体角焊缝）、临时性附件（吊耳、工装等）拆除后的焊痕打磨面，进行 $100\%$ 渗透检测，合格级别不低于 NB/T 47013.5—2015 的 Ⅰ 级。\\[4pt]
\textbf{磁粉检测（MT）}。对铁磁性材料表面和近表面缺陷，可在角焊缝处以磁粉检测替代渗透检测，合格级别不低于 NB/T 47013.4—2015 的 Ⅰ 级。

全部无损检测应在压力试验之前完成且评定合格。检测记录归档保存，随设备出厂资料移交。

\section{压力试验}

\subsection{试验类型与介质}

按 GB/T 150.4—2024 第 4.6 节\cite{gbt150-2024}，选用液压试验（水压试验）作为压力试验方式。试验介质为常温洁净淡水，水质满足 GB 5749—2022 要求\cite{gb5749-2022}。试验前将淡水静置脱气 $60$ min，排除溶解空气。试验环境温度和介质温度均不低于 $5$ ℃，且高于壳体材料韧脆转变温度。



\subsection{试验压力}

按 GB/T 150.4—2024 第 4.6.2 节\cite{gbt150-2024}，液压试验压力由下式确定：
\begin{equation}
P_T = 1.25 P \frac{[\sigma]}{[\sigma]^t}
\label{eq:pt}
\end{equation}

代入：$P_T = 1.25 \times 0.80 \times 189 / 186 = 1.016$ MPa。圆整取 $P_T = 1.02$ MPa。液柱静压力经前文验证可忽略，试验压力以罐顶压力表读数为准。

\subsection{试验应力校核}

试验压力下圆筒环向薄膜应力：
\[
\sigma_T = \frac{P_T (D_i + \delta_e)}{2\delta_e}
         = \frac{1.02 \times 1810.7}{21.4}
         \approx 86.3\ \text{MPa}
\]

液压试验应力校核准则（GB/T 150.4—2024 第 4.6.2.3 条）：
\[
\sigma_T \le 0.9 \phi R_{\text{eL}} = 0.9 \times 0.85 \times 345 = 263.9\ \text{MPa}
\]

$\sigma_T = 86.3$ MPa $ \ll 263.9$ MPa，安全系数 $n_T = 3.06$。封头在试验压力下的最大薄膜应力与筒体相当（$86.3$ MPa），均满足校核条件。

\subsection{试验程序}

试验程序按升压-保压-检查-降压的顺序执行。

升压阶段：分四级升压——$0.20$ MPa（保压 $10$ min，检查初始渗漏）、$0.51$ MPa（$0.5P_T$，保压 $15$ min，检查变形）、$0.80$ MPa（设计压力，保压 $20$ min，锤击检查焊缝）、$1.02$ MPa（试验压力，保压 $30$ min，保压期间不得补压）。

检查阶段：降压至 $0.80$ MPa 的设计压力后全面检查，重点检查纵环焊缝、封头对接焊缝、接管角焊缝、人孔密封面和裙座连接焊缝。检查方法包括目视检查和 $0.5$ kg 圆头小锤轻敲检查（距焊缝中心线 $15$ mm 范围内）。

降压与收尾：检查合格后，以 $\le 0.1$ MPa/min 速率降压至常压，排净罐内积水，干燥压缩空气吹扫内壁至干燥。拆除试压临时封堵和加固装置，恢复设备原状。严禁在保压状态下进行任何修补操作。试验后由施工单位、监理单位和特种设备检验机构三方共同检查试验记录并签署合格报告，试验报告随设备技术档案存档。
